\documentclass[12pt]{article}
\usepackage[margin = 1in]{geometry}

\usepackage{amsmath}
\usepackage{amssymb}
\usepackage{unicode-math}
\usepackage{mleftright}
\usepackage{enumitem}
\usepackage{textcomp, gensymb}
\usepackage{tikz, tikz-cd}
\usepackage[all,cmtip]{xy}

\usepackage{fancyhdr}  % Header and Footer formatting
\pagestyle{fancy}
\renewcommand{\headrulewidth}{0.4pt}
\renewcommand{\footrulewidth}{0.4pt}
\setlength{\headheight}{18pt}

% Header and Footer Information
\lhead{\small\emph{Wenjun Huang}}
\chead{}
\rhead{\textsc{{K-theory seminar}}}
\lfoot{\today}
\cfoot{}
\rfoot{\thepage\ of \ref{NumPages}}  % Counts the pages.

\makeatletter % This provides a total page count as \ref{NumPages}
\AtEndDocument{\immediate\write\@auxout{\string\newlabel{NumPages}{{\thepage}}}}
\makeatother

\usepackage{amsthm}  % This will create the Problem environment
\newtheoremstyle{mythm}%
    {12pt}{}%
    {\sffamily}{}%
    {\bfseries \sffamily}{.}%
    {.5em}%
    {\thmname{#1}\thmnumber{ #2}\thmnote{ #3}}

\theoremstyle{mythm}

\expandafter\let\expandafter\oldproof\csname\string\proof\endcsname
\let\oldendproof\endproof
\renewenvironment{proof}[1][\proofname]{%
  \oldproof[\normalfont \bfseries \sffamily #1]%
}{\oldendproof}

\newtheorem{thm}{Theorem}[section]
\renewcommand*{\proofname}{Proof}
\renewcommand{\qedsymbol}{$\boxed{}$}

\newtheorem{prop}{Propsition}[section]
\renewcommand*{\proofname}{Proof}
\renewcommand{\qedsymbol}{$\boxed{}$}

\newtheorem{cor}{Corollar}[section]
\renewcommand*{\proofname}{Proof}
\renewcommand{\qedsymbol}{$\boxed{}$}

\newtheorem{lemma}{Lemma}[section]
\renewcommand*{\proofname}{Proof}
\renewcommand{\qedsymbol}{$\boxed{}$}

\newtheorem{ex}{Example}[section]
\renewcommand*{\proofname}{Proof}
\renewcommand{\qedsymbol}{$\boxed{}$}

\newtheorem{definition}{Definition}[section]
\renewcommand*{\proofname}{Proof}
\renewcommand{\qedsymbol}{$\boxed{}$}

\newtheorem{claim}{Claim}[section]
\renewcommand*{\proofname}{Proof}
\renewcommand{\qedsymbol}{$\boxed{}$}


\setlist[enumerate]{noitemsep, topsep = 0.2em}
\setlist[enumerate, 1]{label = {\roman*)}}
\setlist[enumerate, 2]{label = {(\alph*)}}
\setlist[description]{topsep = 0.2em, listparindent = \parindent, font = \normalfont}

\newcommand{\ga}{{\symfrak{a}}}
\newcommand{\gb}{{\symfrak{b}}}
\newcommand{\gge}{{\symfrak{e}}}
\newcommand{\gc}{{\symfrak{c}}}
\newcommand{\gp}{{\symfrak{p}}}
\newcommand{\gq}{{\symfrak{q}}}
\newcommand{\gm}{{\symfrak{m}}}
\newcommand{\gn}{{\symfrak{n}}}
\newcommand{\gR}{{\symfrak{R}}}
\newcommand{\bZ}{{\symbb{Z}}}
\newcommand{\bF}{{\symbb{F}}}
\newcommand{\bQ}{{\symbb{Q}}}
\newcommand{\bR}{{\symbb{R}}}
\newcommand{\bC}{{\symbb{C}}}
\newcommand{\id}{{\symrm{id}}}
\newcommand{\CC}{\mathbb{C}}
\newcommand{\R}{\mathbb{R}}
\newcommand{\B}[1]{\boldsymbol{#1}}
\newcommand{\mC}{\mathcal{C}}
\newcommand{\dsp}{\displaystyle}
\newcommand{\Q}{\mathbb{Q}}
\newcommand{\N}{\mathbb{N}}
\newcommand{\RR}{\mathbb{R}}
\newcommand{\ZZ}{\mathbb{Z}}
\newcommand{\A}{\mathcal A}
\newcommand{\eps}{\varepsilon}
\newcommand*{\defeq}{\mathrel{\vcenter{\baselineskip0.5ex \lineskiplimit0pt
                     \hbox{\scriptsize.}\hbox{\scriptsize.}}}%
     		     =}
\DeclareMathOperator{\im}{im}
\DeclareMathOperator{\Ann}{Ann}
\DeclareMathOperator{\Spec}{Spec}
\DeclareMathOperator{\Max}{Max}
\DeclareMathOperator{\Supp}{Supp}
\DeclareMathOperator{\Ind}{Ind}
\newcommand{\tprime}{\ensuremath{\prime}}

\begin{document}
\section{Cylindrical algebraic decomposition}
\begin{thm}[Thom lemma]
    Let \( f_1, \ldots, f_k \in \mathbb{R}[T] \) be nonzero polynomials such that if \( f_i' \neq 0 \), then \( f_i' \in \{f_1, \ldots, f_k\} \). Let \(\epsilon : \{1, \ldots, k\} \to \{-1, 0, 1\}\), and put

\[
A_\epsilon := \{t \in \mathbb{R} : \mathrm{sign}(f_i(t)) = \epsilon(i), \, i = 1, \ldots, k\}, \text{ a subset of } \mathbb{R}.
\]

Then \( A_\epsilon \) is empty, a point, or an interval.

Moreover,  If \( A_\epsilon \neq \emptyset \), then its closure is given by

\[
\mathrm{cl}(A_\epsilon) = \{t \in \mathbb{R} : \mathrm{sign}(f_i(t)) \in \{\epsilon(i), 0\}, \, i = 1, \ldots, k\}.
\]

If \( A_\epsilon = \emptyset \),then \(cl(A_\epsilon)\) is empty or points.
\end{thm}

\begin{proof}
   By induction on \(s\). The Lemma holds trivially for \(s=0\). Let \(f_{1},\ldots,f_{s},f_{s+1}\in R[\mathbf{x}]\setminus\{0\}\) be polynomials such that if \(f^{\prime}_{k}\neq 0\), then \(f^{\prime}_{k}\in\{f_{1},\ldots,f_{s+1}\}\). Without loss of generality we assume that \(\deg(f_{s+1})=\max\{\deg(f_{k}):1\leqslant k\leqslant s+1\}\).

Let \(\varepsilon^{\prime}\colon\{1,\ldots,s,s+1\}\to\{-1,0,1\}\) and \(\varepsilon\colon\{1,\ldots,s\}\to\{-1,0,1\}\) the restriction.

Note that
\[
A_{\varepsilon^{\prime}}=A_{\varepsilon}\cap\{x\in R:\operatorname{sign}(f_{s+1}(x))=\varepsilon^{\prime}(s+1)\}.
\]

By induction \(A_{\varepsilon}\) is empty, a point, or an interval.

If \(A_{\varepsilon}\) is empty or a point, then obviously so is \(A_{\varepsilon^{\prime}}\) and the other property follows immediately by induction hypothesis on \(A_{\varepsilon}\).

Assume \(A_{\varepsilon}\) is an interval. Now \(f^{\prime}_{s+1}=0\) or \(f^{\prime}_{s+1}\in\{f_{1},\ldots,f_{s}\}\). So by definition of \(A_{\varepsilon}\), \(f^{\prime}_{s+1}\) has constant sign on \(A_{\varepsilon}\). Therefore \(f_{s+1}\) is either strictly increasing, or strictly decreasing or constant on \(A_{\varepsilon}\).

Consider \(A_{\varepsilon}=(a,b)\). There are three cases depending on \(\varepsilon^{\prime}(s+1)\):

\textbf{Case 1.} \(A_{\varepsilon^{\prime}}=\{x\in(a,b):f_{s+1}(x)>0\}\).

\textbf{Case 2.} \(A_{\varepsilon^{\prime}}=\{x\in(a,b):f_{s+1}(x)<0\}\).

\textbf{Case 3.} \(A_{\varepsilon^{\prime}}=\{x\in(a,b):f_{s+1}(x)=0\}\).

If \(A_{\varepsilon^{\prime}}=\emptyset\) there is nothing to prove.

Assume \(A_{\varepsilon^{\prime}}\neq\emptyset\). If \(f_{s+1}\) is constant on \(A_{\varepsilon}\) then \(f_{s+1}\) is a constant polynomial \(f_{s+1}(x)=c\neq 0\). So \(A_{\varepsilon^{\prime}}\) is empty or \(A_{\varepsilon^{\prime}}=(a,b)\) depending on whether \(\operatorname{sign}(c)=\varepsilon^{\prime}(s+1)\).

Assume now \(f_{s+1}\) strictly increasing on \(A_{\varepsilon}\) and \(A_{\varepsilon^{\prime}}=\{x\in(a,b):f_{s+1}(x)>0\}\neq\emptyset\). Let \(x_{0}=\inf\{x\in(a,b):f_{s+1}(x)>0\}\). Since \(f_{s+1}\) is strictly increasing it follows that \(f_{s+1}(x)>0\) \(\forall\,x\in(a,b)\) with \(x>x_{0}\). So \(A_{\varepsilon^{\prime}}=(x_{0},b)\) and its closure is \([x_{0},b]=A_{\varepsilon^{\prime}}\). The other cases are treated similarly.

\end{proof}



\textbf{Notation.}
$X$ is a nonempty topological space, $E$ a ring of continuous real-valued functions
$f: X \to \mathbb{R}$, the ring operations being pointwise addition and multiplication, with
multiplicative identity the function on $X$ that takes the constant value $1$.

For $f(T)$ as above and $z \in X$, we write $f(z, T)$ for the polynomial

\[
f(z,T) = f_0(z) + f_1(z)T + \cdots + f_d(z)T^d \in \mathbb{R}[T].
\]

Let us say that a polynomial $f(T) \in E[T]$ has \emph{zero-free leading coefficient} if  
\[
f(T) = c_d T^d + c_{d-1} T^{d-1} + \cdots + c_0
\]  
with $c_i \in E$, such that $c_d$ has no zero on $X$.
\begin{definition}
Call a set \( A \subseteq X \) an \( E \)-set if \( A \) is a finite union of sets of the form  
\[\{x \in X : f(x) = 0, \, g_1(x) > 0, \ldots, g_k(x) > 0\}, \text{ with } f, g_1, \ldots, g_k \in E.\]  

\end{definition}
Note that the \( E \)-sets form a boolean algebra of subsets of \( X \).unions, intersections and
 taking complements.

\begin{definition}
    If \( X = \mathbb{R}^m \), and \( E = \mathbb{R}[x_1, \ldots, x_m] \) is the ring of real polynomial functions on \( X \), then the \( E \)-sets are exactly the semialgebraic subsets of \( \mathbb{R}^m \).
\end{definition}

\begin{lemma}[Lemma 2]
Let \( G(A, T) = A_0 + A_1T + \cdots + A_dT^d \in \ZZ[A, T], A = (A_0, \cdots, A_d) \), be a general polynomial of degree \( d \), and \( e \in \{0, \cdots, d, \infty\} \). Then the set
\[
\{a \in \RR^{d+1} : G(a, T) \text{ has exactly } e \text{ distinct complex zeros}\}
\]
is a finite union of sets of the form
\[
\{a \in \CC^{d+1} : p_1(a) = \cdots = p_k(a) = 0, \, q(a) \neq 0\},
\]
where \( p_i, q \in \ZZ[A] \).
\end{lemma}

\begin{proof}
The cases \( d = 0 \) or \( e \in \{0, \infty\} \) are trivial.  
Let \( d > 0, e \in \{1, \cdots, d\} \), and \( a = (a_0, \cdots, a_d) \in \CC^{d+1}, a_d \neq 0 \).  

Let \( g = \) degree of GCD\((G(a, T), \frac{\partial G}{\partial T}(a, T))\) in \(\CC[T]\).  

Then the number of distinct complex zeros of \( G(a, T) \) is \( d - g \), and the degree of LCM\((G(a, T), \frac{\partial G}{\partial T}(a, T))\) is \( 2d - g - 1 \).  

Hence the condition is that \( G(a, T) \) has at most \( e \) distinct zeros, which is equivalent to \( d - g \leq e \), that is, to \( 2d - g - 1 \leq d + e - 1 \).  

The last condition is equivalent to the condition:  
(*) There exist \( q(x, T) = x_0 + x_1 T + \cdots + x_{e-1} T^{e-1} \) and  
\[r(x, T) = x_e + x_{e+1} T + \cdots + x_{2e} T^e,\]  
with  
\[x = (x_0, \cdots, x_{2e}) \in \CC^{2e+1} \setminus 0,\]  
such that  
\[G(a, T)q(x, T) = \frac{\partial G}{\partial T}(a, T) r(x, T).\]  
This equality can be rewritten as  
\[G(a, T)q(x, T) - \frac{\partial G}{\partial T}(a, T) r(x, T) = \beta_0(a, x) + \beta_1(a, x) T + \cdots + \beta_{d+e-1}(a, x) T^{d+e-1}\]  
where \( \beta = (\beta_0, \cdots, \beta_{d+e-1}) : \CC^{d+1} \times \CC^{2e+1} \to \CC^{d+e-1} \) is a bilinear function.

So (*) is equivalent to the condition \(\beta_0(a, x) = \cdots = \beta_{d+e-1}(a, x) = 0\) has nonzero solution \(x \in \CC^{2e+1}\).  
The last condition is equivalent to the vanishing of all \((2e+1)\)-minor of the matrix of the linear map \(\beta(a, \cdot)\).  
Note that each of the minors is a polynomial in \(a_0, \cdots, a_d\) with coefficients in \(\ZZ\).  
Therefore, for each \(d' \leq d\), the set \(M_e^{d'} = \{a \in \RR^{d+1}: G(a, T) \text{ is of degree } d' \text{ and has at most } e \text{ distinct complex zeros}\}\) is the intersection of the set \(\{a \in \RR^{d+1}: a_d = \cdots = a_{d'+1} = 0, a_{d'} \neq 0\}\) with the zero set of certain polynomials in \(\ZZ[A]\).  
So \(\{a = (a_0, \cdots, a_d) \in \RR^{d+1}: G(a, T) \text{ has exactly } e \text{ complex zeros}\} = \bigcup_{d'=0}^d M_e^{d'} \setminus M_{e-1}^{d'} \) a semi-algebraic set.
\end{proof}

\begin{cor}
As a consequence, for every \( f \in \RR[X_1, \cdots, X_n, T] = \RR[X][T] \),
\[f(X_1, \cdots, X_n, T) = a_0(X) + a_1(X)T + \cdots + a_d(X)T^d,\]
the set
\[
\{x \in \RR^n : f(x, T) \text{ has exactly } e \text{ distinct complex zeros}\}
\]
is a semi-algebraic subset of \(\RR^n\).
\end{cor}

\begin{proof}
Since \(f(x, T) = G(a_0(x), \cdots, a_d(x), T)\), the corollary follows from Lemma 2.
\end{proof}




\begin{thm}[Cylindrical algebraic decomposition]
    Let \( f_1(T), \ldots, f_M(T) \in E[T] \). Then the list \( f_1, \ldots, f_M \) can be augmented to a list \( f_1, \ldots, f_N \) in \( E[T] \) (\( M \leq N \)), and \( X \) can be partitioned into finitely many \( E \)-sets \( X_i \) (\( 1 \leq i \leq k \)) such that for each connected component \( C \) of each \( X_i \) there are continuous real-valued functions \( \xi_{C,1} < \cdots < \xi_{C,m(C)} \) on \( C \) with the following two properties:

\begin{enumerate}
    \item Each function \( f_n \) (\( 1 \leq n \leq N \)) has constant sign \( (-1, 0, \text{or } 1) \) on each of the graphs \( \Gamma(\xi_{C,j}) \) (\( 1 \leq j \leq m(C) \)) and on each of the sets \( (\xi_{C,j}, \xi_{C,j+1}) \) (\( 0 \leq j \leq m(C) \)), where \( \xi_{C,0} := -\infty \) and \( \xi_{C,m(C)+1} := +\infty \) are constant functions on \( C \) by convention;

    \item Each of the sets \( \Gamma(\xi_{C,j}) \) and \( (\xi_{C,j}, \xi_{C,j+1}) \) from (1) is of the form
    \[
    \{(x,t) \in C \times \mathbb{R}: \mathrm{sign}(f_n(x,t)) = \epsilon(n) \text{ for } n = 1, \ldots, N\}
    \]
    for a suitable sign condition \( \epsilon: \{1, \ldots, N\} \rightarrow \{-1,0,1\} \).
\end{enumerate}
\end{thm}


\begin{proof}

Let \( d = \max\{\deg_T(f_i), i = 1, \cdots, p\} \).  
Augment \( f_1, \cdots, f_p \) to  

\[
\{f_1, \cdots, f_{p+q}\} = \left\{\frac{\partial^\nu f_i}{\partial T^\nu} : 1 \leq i \leq p, 0 \leq \nu \leq d\right\}.
\]

For each \(\Delta \subset \{1, \cdots, p\} \times \{0, \cdots, d\}\), and \( e \in \{0, \cdots, pd^2\} \cup \{\infty\} \), put  

\[
f_\Delta(T) = \prod_{(i, \nu) \in \Delta} \frac{\partial^\nu f_i}{\partial T^\nu} \in \mathbb{R}[X][T], \text{ and }
\]

\[
A_{\Delta,e} = \{x \in \mathbb{R}^n : f_\Delta(x,T) \text{ has exactly } e \text{ complex zeros}\}.
\]

By Lemma 2, \( A_{\Delta,e} \) is a semi-algebraic set.  
For a given \(\Delta\) the family \(\{A_{\Delta,e} : e \text{ varies }\}\) forms a partition of \(\mathbb{R}^n\). Since the class of semi-algebraic sets is a boolean algebra we can find a partition (the intersection of the partitions) \(\mathbb{R}^n = S_1 \cup \cdots \cup S_k\), where each \(S_i\) is semi-algebraic such that each set \(A_{\Delta,e}\) is a union of the \(S_i's\). We will prove that \(f_1, \cdots, f_{p+q}\) and \(S_1, \cdots, S_k\) satisfy the conclusion of the theorem.

For each connected component \( C \) of \( S_i \) put

\[
\Delta(C) = \left\{(i, \nu) : \frac{\partial^\nu f_i}{\partial T^\nu} \neq 0 \text{ on } C \times \mathbb{R}\right\}.
\]

By Lemma 3, there exist continuous functions \(\xi_{C,1} < \cdots < \xi_{C,r}(C)\) on \( C \) such that \(\{(x,t) \in C \times \mathbb{R} : f_\Delta(C) = 0\} = \Gamma(\xi_{C,1}) \cup \cdots \cup \Gamma(\xi_{C,r}(C))\).

\textbf{Check (i):} If \((i, \nu) \notin \Delta(C)\) then \(\frac{\partial^\nu f_i}{\partial T^\nu} \equiv 0\) on the sets given in (i).

If \((i, \nu) \in \Delta(C)\), then \(C \subset A_{\{(i,\nu)\},e}\), for certain \(e \in \{0, \cdots, d\} \cup \{\infty\}\) and the number of real zeros of \(\frac{\partial^\nu f_i}{\partial T^\nu}(x,T)\) is independent of \(x \in C\).

Since \(\frac{\partial^\nu f_i}{\partial T^\nu}\) is a factor of \(f_\Delta(C)\), by Lemma 3, the zeros of \(\frac{\partial^\nu f_i}{\partial T^\nu}(x,T)\), for \(x \in C\), must be among the \(\xi_{C,j}(x)'s\). Since \(C\) is connected, (i) is checked.

\textbf{Check (ii):} Let \( B \) be one of the sets in (i). By (i), \( \epsilon(i, \nu) = \text{sign}(\frac{\partial^\nu f_i}{\partial T^\nu} |_B) \) is well-defined. Put

\[
B' = \left\{(x, t) \in C \times \mathbb{R}: \text{sign}\left(\frac{\partial^\nu f_i}{\partial T^\nu}(x, t)\right) = \epsilon(i, \nu), \, 1 \leq i \leq p, 0 \leq \nu \leq d\right\}.
\]

Clearly \( B \subset B' \). If \( B \neq B' \) then exist \((x, t') \in B' \setminus B\), \((x, t) \in B\) (say \( t < t'\)). Thom's lemma implies that \(\{r \in \mathbb{R} : (x, r) \in B'\}\) is connected, so \(\{x\} \times [t, t'] \subset B'\). Since \((x, t) \in B\), \((x, t') \notin B\), \(f_\Delta(c)\) must change sign on \(\{x\} \times [t, t']\). But \(f_\Delta(c)\) is a product of \(\frac{\partial^\nu f_i}{\partial T^\nu}\), so \(f_\Delta(c)\) cannot change sign on \(B'\), contradiction. Therefore \(B = B'\).




\end{proof}

\begin{thm}[Lojasiewicz]
The number of connected components of a semi-algebraic set is finite,
and each of the components is also semi-algebraic.
\end{thm}


\begin{cor}
     If $(X, E)$ has the Lojasiewicz property, then $(X \times \mathbb{R}, E(T))$ also has the Lojasiewicz property, and the image of any $E[T]$-set $S \subseteq X \times \mathbb{R}$ under the projection map $X \times \mathbb{R} \to X$ is an $E$-set.
\end{cor}

\begin{proof}
 
Write an \( E[T] \)-set \( S \) as a finite union of sets of the form
\[
\{(x,t) \in X \times \mathbb{R} : f(x,t) = 0, \, g_1(x,t) > 0, \ldots, g_k(x,t) > 0\}
\]
with \( f, g_1, \ldots, g_k \in E[T] \). Let \( f_1, \ldots, f_M \) be a list of all the polynomials in \( E[T] \) that are involved in such a definition of \( S \). Now apply the theorem to this list.
\end{proof}




\begin{thm}[Tarski-Seidenberg Theorem]
    The $\R[T_1, \ldots, T_m]$-sets in the space $\R^m$ are exactly the semialgebraic subsets of $\R^m$. For $m = 0$ the space $\RR^m$ has only one point, so the pair $(\RR^m, \RR[T_1, \ldots, T_m])$ obviously has the Lojasiewicz property for $m = 0$; hence by induction on $m$, using (2.9), it follows that $(\RR^m, \RR[T_1, \ldots, T_m])$ has the Lojasiewicz property for each $m$, and that the image of a semialgebraic subset of $\RR^{m+1}$ under the projection map $\RR^{m+1} \to \RR^m$ is a semialgebraic subset of $\RR^m$.
\end{thm}
We can now draw the following conclusion about the model-theoretic structure $(\R, <, 0, 1, +, -, \cdot)$, which we will call ``the ordered field of real numbers''.

\begin{cor}
The sets $S \subseteq \R^m$ ($m = 0, 1, 2, \ldots$) that are definable using constants in the ordered field of real numbers are exactly the semialgebraic sets.The ordered field of real numbers is o-minimal.
\end{cor}
\begin{proof}
   
 semialgebraic sets form a structure on \((\mathbb{R}, <)\). Clearly the primitives of the ordered field of reals are semialgebraic, the semialgebraic sets are definable using constants in the ordered field of real numbers. Finally, the semialgebraic subsets of \(\mathbb{R}\) are clearly the finite unions of intervals and points. The desired result follows.



\end{proof}


\section{}
\subsection{noether normalization}

Let \( f(X_1, \ldots, X_n) \) be a nonzero polynomial of degree \( d \) over a field \( K \) of characteristic 0, say 
\[
f = \sum_{|i|=d} c_i X^i + \text{terms of lower degree}, \quad n > 0.
\]
Here \( i = (i_1, \ldots, i_n) \in \mathbb{N}^n, |i| = i_1 + \cdots + i_n \) and \( X^i = X_1^{i_1} \cdots X_n^{i_n}, c_i \in K \). Let \( (\lambda_1, \ldots, \lambda_{n-1}) \in K^{n-1} \). We are going to make the (invertible) substitution

\begin{align*}
X_1 &\to X_1 + \lambda_1 X_n, \\
&\vdots \\
X_{n-1} &\to X_{n-1} + \lambda_{n-1} X_n, \\
X_n &\to X_n.
\end{align*}

Note that, with \(\lambda^i = \lambda_1^{i_1} \cdots \lambda_{n-1}^{i_{n-1}}\), this substitution transforms \(f\) into

\begin{align*}
&f(X_1 + \lambda_1 X_n, \ldots, X_{n-1} + \lambda_{n-1} X_n, X_n) \\
&= \left( \sum_{|i|=d} c_i \lambda^i \right) X_n^d + \text{terms of degree } < d \text{ in } X_n.
\end{align*}

Since the polynomial \(\sum_{|i|=d} c_i Y_1^{i_1} \cdots Y_{n-1}^{i_{n-1}} \in K[Y_1, \ldots, Y_{n-1}]\) is nonzero, it follows that for ``almost all'' \(\lambda = (\lambda_1, \ldots, \lambda_{n-1}) \in K^{n-1}\), namely for all those \(\lambda\) that are not zeros of this polynomial, the transformed polynomial

\[
f(X_1 + \lambda_1 X_n, \ldots, X_{n-1} + \lambda_{n-1} X_n, X_n)
\]

is of (total) degree \(d\), and has a term of the form \(cX_n^d\) with \(c \in K - \{0\}\).
\subsection{}
\begin{definition}

Let \( f_1(T), \ldots, f_M(T) \in E[T] \), and let \( C \) be a connected subset of \( X \). A \emph{decomposition of \( f_1, \ldots, f_M \) above \( C \)} consists of continuous real-valued functions \( \xi_1 < \cdots < \xi_r \) on \( C \) such that

\begin{enumerate}
    \item each function \( f_m \) has constant sign on each set \( \Gamma(\xi_i) \) (\( 1 \leq i \leq r \)) and on each set \( (\xi_i, \xi_{i+1}) \) (\( 0 \leq i \leq r \)), where \( \xi_0 := -\infty, \xi_{r+1} := +\infty \);
    
    \item each set \( \Gamma(\xi_i) \) is contained in the zero set of some function \( f_m \).
\end{enumerate}
    
\end{definition}


Note: if (1) holds, then (2) can be satisfied by removing the $\xi_i$'s not satisfying (2).  

Let us say that a polynomial $f(T) \in E[T]$ has \emph{zero-free leading coefficient} if  
\[
f(T) = c_d T^d + c_{d-1} T^{d-1} + \cdots + c_0
\]  
with $c_i \in E$, such that $c_d$ has no zero on $X$.

if $f(T) \in E[T]$ has zero-free leading coefficient and a point $y \in X$ is given, then the real zeros of $f(x, T)$ are bounded in absolute value by some constant, for all $x$ in a suitably small neighborhood of $y$.

\begin{lemma}
Suppose the polynomials \( f_1, \ldots, f_N \in E[T] \) have zero-free leading coefficients, and each nonzero partial \(\partial f_n / \partial T\) also belongs to \(\{f_1, \ldots, f_N\}\). Let  
\[\xi_1 < \cdots < \xi_r\]  
be a decomposition of \( f_1, \ldots, f_N \) above a connected set \( C \subseteq X \). Fix  
\[i \in \{1, \ldots, r\}, \text{ let } \xi = \xi_i, \text{ put } \epsilon(n) = \mathrm{sign}(f_n | \Gamma(\xi)) \text{ for } 1 \leq n \leq N.\]  
Then  

\begin{enumerate}
    \item[(a)] 
    \[\Gamma(\xi) = \{(x, t) \in C \times \mathbb{R} : \mathrm{sign}(f_n(x, t)) = \epsilon(n) \text{ for } n = 1, \ldots, N\},\]
    
    \item[(b)] 
    \[\xi : C \to \mathbb{R} \] extends uniquely to a continuous function \[\eta : \mathrm{cl}(C) \to \mathbb{R}, \]eand \[ \mathrm{cl}(\Gamma(\xi)) = \Gamma(\eta) = \{(y, t) \in \mathrm{cl}(C) \times \mathbb{R} : \mathrm{sign}(f_n(y, t)) \in \{\epsilon(n), 0\} \text{ for } n = 1, \ldots, N\}.\]
\end{enumerate}
\end{lemma}
\begin{proof}
     For fixed \( x \in C \) the set  
\[A_{x,\epsilon} := \{ t \in \mathbb{R} : \text{ sign } (f_n(x,t)) = \epsilon(n) \text{ for } n = 1, \ldots, N\}\]  
contains \(\{\xi(x)\}\) and is finite since some \( f_m \) vanishes on \(\Gamma(\xi)\) and \( f_m(x,T) \neq 0 \). The polynomials \( f_1(x,T), \ldots, f_N(x,T) \) satisfy the hypothesis of Thom's lemma, so \( A_{x,\epsilon} \) is connected. Hence \( A_{x,\epsilon} = \{\xi(x)\} \). This proves (a).
For (b), take \( m \) such that \( f_m(x, \xi(x)) = 0 \) for all \( x \in C \). Let \( y \in cl(C) \). Since \( f_m(T) \) has zero-free leading coefficient, there is a neighborhood \( U \) of \( y \) such that \( \xi \) is bounded on \( U \cap C \). Hence the limits

\[l := \liminf_{x \to y} \xi(x)\]

and \( L := \limsup_{x \to y} \xi(x) \)

are real numbers (not \(\pm \infty\)). Now \( sign(f_n(x, \xi(x))) = \epsilon(n) \) for all \( n \), hence \( sign(f_n(y, l)) \in \{\epsilon(n), 0\} \), and \( sign(f_n(y, L)) \in \{\epsilon(n), 0\} \) for all \( n \). By Thom's lemma the set \(\{t \in \mathbb{R}: sign(f_n(y, t)) \in \{\epsilon(n), 0\} \text{ for } n = 1, \ldots, N\}\) is connected, so this set contains \([l, L]\). With \( m \) as above we have \( f_m(y, t) = 0 \) for all \( t \in [l, L] \), and because \( f_m(y, T) \neq 0 \), this implies \( l = L \). Since \( y \) was an arbitrary point of \( cl(C) \) we have shown that \( \eta(y) := \lim_{x \to y} \xi(x) \) defines a function \( \eta: cl(C) \to \mathbb{R} \). Continuity of \( \xi \) implies continuity of \( \eta \), so \( \Gamma(\eta) \) is closed, hence \( cl(\Gamma(\xi)) = \Gamma(\eta) \).

The argument with Thom's lemma that we just gave also shows

\[\Gamma(\eta) = \{(y, t) \in cl(C) \times \mathbb{R}: sign(f_n(y, t)) \in \{\epsilon(n), 0\} \text{ for } n = 1, \ldots, N\}.\]

This finishes the proof of (b).
\end{proof}

\begin{lemma}
With the same assumptions as in the previous lemma, set \(\xi_0 := -\infty\),  
\(\xi_{r+1} := +\infty\). Fix an \(i \in \{0, \ldots, r\}\) and put \(\epsilon(n) = \mathrm{sign}(f_n|(\xi_i, \xi_{i+1}))\). Then we have  

\begin{enumerate}
    \item[(a)] \((\xi_i, \xi_{i+1}) = \{(x, t) \in C \times \mathbb{R}: \mathrm{sign}(f_n(x, t)) = \epsilon(n) \text{ for } n = 1, \ldots, N\}\),
    
    \item[(b)] \(\mathrm{cl}(\xi_i, \xi_{i+1}) = \{(y, t) \in \mathrm{cl}(C) \times \mathbb{R}: \mathrm{sign}(f_n(y, t)) \in \{\epsilon(n), 0\} \text{ for } n = 1, \ldots, N\}\).
\end{enumerate}
\end{lemma}
\begin{proof}
    ### Step 1. Understanding the Context

We assume the setting from a previous lemma:  
- \( C \) is a connected set (e.g., a cell in a cylindrical algebraic decomposition).  
- \( f_1, \ldots, f_N \) are continuous functions on \( C \times \mathbb{R} \).  
- There exist continuous functions \( \xi_1, \ldots, \xi_r : C \to \mathbb{R} \) such that for each \( x \in C \), the real line is partitioned into intervals  
  \[
  (-\infty, \xi_1(x)),\ (\xi_1(x), \xi_2(x)),\ \ldots,\ (\xi_r(x), +\infty)
  \]  
  on which the signs of the \( f_n \) are constant.  
- Define \( \xi_0(x) = -\infty \), \( \xi_{r+1}(x) = +\infty \).  
- For a fixed \( i \in \{0, \ldots, r\} \), define the set  
  \[
  (\xi_i, \xi_{i+1}) = \left\{ (x, t) \in C \times \mathbb{R} : \xi_i(x) < t < \xi_{i+1}(x) \right\}.
  \]  
- On this set, the sign of each \( f_n \) is constant; denote it by \( \epsilon(n) \).  
- The closure \( \mathrm{cl}(C) \) is taken in the ambient space, and the functions \( \xi_i \) extend continuously to \( \mathrm{cl}(C) \).

---

### Step 2. Statement of the Lemma

With the above assumptions, define \( \epsilon(n) = \mathrm{sign}(f_n|_{(\xi_i, \xi_{i+1})}) \). Then:

1. **(a)**  
   \[
   (\xi_i, \xi_{i+1}) = \left\{ (x, t) \in C \times \mathbb{R} : \mathrm{sign}(f_n(x, t)) = \epsilon(n) \text{ for } n = 1, \ldots, N \right\}.
   \]

2. **(b)**  
   \[
   \mathrm{cl}(\xi_i, \xi_{i+1}) = \left\{ (y, t) \in \mathrm{cl}(C) \times \mathbb{R} : \mathrm{sign}(f_n(y, t)) \in \{\epsilon(n), 0\} \text{ for } n = 1, \ldots, N \right\}.
   \]

---

### Step 3. Proof of Part (a)

Let  
\[
A = (\xi_i, \xi_{i+1}), \quad S = \left\{ (x, t) \in C \times \mathbb{R} : \mathrm{sign}(f_n(x, t)) = \epsilon(n) \text{ for all } n \right\}.
\]

- **Show \( A \subseteq S \):**  
  By definition, for all \( (x, t) \in A \), we have \( \mathrm{sign}(f_n(x, t)) = \epsilon(n) \). Hence, \( A \subseteq S \).

- **Show \( S \subseteq A \):**  
  Let \( (x, t) \in S \). Then \( \mathrm{sign}(f_n(x, t)) = \epsilon(n) \) for all \( n \), so no \( f_n(x, t) \) is zero.  
  By the previous lemma, \( t \) must lie in one of the intervals \( (\xi_j(x), \xi_{j+1}(x)) \).  
  If \( j \ne i \), then the sign pattern on \( (\xi_j(x), \xi_{j+1}(x)) \) would be different from \( \epsilon(n) \), a contradiction.  
  Hence, \( j = i \), so \( (x, t) \in A \). Thus, \( S \subseteq A \).

✅ Therefore, \( A = S \).

---

### Step 4. Proof of Part (b)

Let  
\[
A = (\xi_i, \xi_{i+1}), \quad B = \left\{ (y, t) \in \mathrm{cl}(C) \times \mathbb{R} : \mathrm{sign}(f_n(y, t)) \in \{\epsilon(n), 0\} \text{ for all } n \right\}.
\]

We prove \( \mathrm{cl}(A) = B \).

#### (i) Show \( \mathrm{cl}(A) \subseteq B \)

Let \( (y, t) \in \mathrm{cl}(A) \). Then there exists a sequence \( \{(x_k, t_k)\} \subset A \) such that \( (x_k, t_k) \to (y, t) \).  
Since \( A \subset C \times \mathbb{R} \), we have \( y \in \mathrm{cl}(C) \).  
For each \( n \), since \( (x_k, t_k) \in A \), we have \( \mathrm{sign}(f_n(x_k, t_k)) = \epsilon(n) \).  
By continuity of \( f_n \), \( f_n(x_k, t_k) \to f_n(y, t) \).  
- If \( f_n(y, t) \ne 0 \), then for large \( k \), \( \mathrm{sign}(f_n(x_k, t_k)) = \mathrm{sign}(f_n(y, t)) \), so \( \mathrm{sign}(f_n(y, t)) = \epsilon(n) \).  
- If \( f_n(y, t) = 0 \), then \( \mathrm{sign}(f_n(y, t)) = 0 \).  

In both cases, \( \mathrm{sign}(f_n(y, t)) \in \{\epsilon(n), 0\} \). Hence, \( (y, t) \in B \).  
✅ So, \( \mathrm{cl}(A) \subseteq B \).

#### (ii) Show \( B \subseteq \mathrm{cl}(A) \)

Let \( (y, t) \in B \). Then \( y \in \mathrm{cl}(C) \), and for all \( n \), \( \mathrm{sign}(f_n(y, t)) \in \{\epsilon(n), 0\} \).  
We construct a sequence in \( A \) converging to \( (y, t) \).

- Since \( y \in \mathrm{cl}(C) \), there exists \( \{x_k\} \subset C \) such that \( x_k \to y \).  
- By continuity of \( \xi_i, \xi_{i+1} \) on \( \mathrm{cl}(C) \), we have  
  \[
  \xi_i(x_k) \to \xi_i(y), \quad \xi_{i+1}(x_k) \to \xi_{i+1}(y).
  \]

We now choose \( t_k \in \mathbb{R} \) such that \( \xi_i(x_k) < t_k < \xi_{i+1}(x_k) \) and \( t_k \to t \).  
This ensures \( (x_k, t_k) \in A \). Consider cases:

- **Case 1:** \( \xi_i(y) < t < \xi_{i+1}(y) \)  
  Then for large \( k \), \( \xi_i(x_k) < t < \xi_{i+1}(x_k) \), so take \( t_k = t \).

- **Case 2:** \( t = \xi_i(y) \) and \( \xi_i(y) < \xi_{i+1}(y) \)  
  Choose \( t_k = \xi_i(y) + \delta_k \), where \( \delta_k \to 0^+ \) and \( \xi_i(x_k) < t_k \).  
  This is possible since \( \xi_i(x_k) \to \xi_i(y) \).  
  Also, for large \( k \), \( t_k < \xi_{i+1}(x_k) \) since \( \xi_{i+1}(x_k) \to \xi_{i+1}(y) > \xi_i(y) \).

- **Case 3:** \( t = \xi_{i+1}(y) \) and \( \xi_i(y) < \xi_{i+1}(y) \)  
  Choose \( t_k = \xi_{i+1}(y) - \delta_k \), with \( \delta_k \to 0^+ \), so that \( t_k > \xi_i(x_k) \) and \( t_k < \xi_{i+1}(x_k) \).

- **Case 4:** \( \xi_i(y) = \xi_{i+1}(y) \)  
  Then necessarily \( t = \xi_i(y) \). Define  
  \[
  t_k = \frac{\xi_i(x_k) + \xi_{i+1}(x_k)}{2}.
  \]  
  Then \( \xi_i(x_k) < t_k < \xi_{i+1}(x_k) \), and \( t_k \to \frac{\xi_i(y) + \xi_{i+1}(y)}{2} = \xi_i(y) = t \).

In all cases, we have \( (x_k, t_k) \in A \) and \( (x_k, t_k) \to (y, t) \). Hence, \( (y, t) \in \mathrm{cl}(A) \).  
✅ So, \( B \subseteq \mathrm{cl}(A) \).

---

### Step 5. Conclusion

We have shown:
- \( (\xi_i, \xi_{i+1}) = \left\{ (x, t) \in C \times \mathbb{R} : \mathrm{sign}(f_n(x, t)) = \epsilon(n) \text{ for all } n \right\} \),
- \( \mathrm{cl}(\xi_i, \xi_{i+1}) = \left\{ (y, t) \in \mathrm{cl}(C) \times \mathbb{R} : \mathrm{sign}(f_n(y, t)) \in \{\epsilon(n), 0\} \text{ for all } n \right\} \).

This completes the proof of the lemma.

---

### Final Answer

\[
\boxed{
\begin{array}{l}
\text{(a) } (\xi_i, \xi_{i+1}) = \left\{ (x, t) \in C \times \mathbb{R} : \mathrm{sign}(f_n(x, t)) = \epsilon(n) \text{ for } n = 1, \ldots, N \right\}, \\
\text{(b) } \mathrm{cl}(\xi_i, \xi_{i+1}) = \left\{ (y, t) \in \mathrm{cl}(C) \times \mathbb{R} : \mathrm{sign}(f_n(y, t)) \in \{\epsilon(n), 0\} \text{ for } n = 1, \ldots, N \right\}.
\end{array}
}
\]
\end{proof}

\begin{definition}
Let \( F \) be a finite set of continuous real-valued functions on a topological space \( X \). We say that a set \( S \subseteq X \) is \textbf{described by} \( F \) if \( S \) is a finite union of sets of the form  
\[e(F) := \{x \in X : \text{ sign}(f(x)) = e(f) \text{ for all } f \in F\}, \text{ with } e: F \to \{-1, 0, 1\}.\]  
Since \( F \) is a finite collection of functions, the sets described by \( F \) form a finite boolean algebra of subsets of \( X \).
\end{definition}

\begin{cor}
    A closed semialgebraic subset of $\mathbb{R}^n$ is a finite union of sets of the form 
\[\{x \in \mathbb{R}^n : f_1(x) > 0, \ldots, f_k(x) > 0\}\]
with $f_1, \ldots, f_k \in \mathbb{R}[X_1, \ldots, X_n]$.
\end{cor}
\begin{proof}
    ### Step 1. Understanding the Corollary

The corollary states:

> A **closed** semialgebraic subset of \(\mathbb{R}^n\) is a finite union of sets of the form  
> \[
\{x \in \mathbb{R}^n : f_1(x) > 0, \ldots, f_k(x) > 0\}
\]  
> with \(f_1, \ldots, f_k \in \mathbb{R}[X_1, \ldots, X_n]\).

However, sets defined by strict inequalities \(f_i(x) > 0\) are **open** in the Euclidean topology (since polynomials are continuous). A finite union of open sets is open, so such a representation would imply that every closed semialgebraic set is also open — which is false in general (e.g., a closed interval in \(\mathbb{R}\) is not open).

On the other hand, it is a standard fact that **open** semialgebraic sets admit such a representation. Therefore, we interpret the corollary as referring to **open** semialgebraic sets.

---

### Step 2. Preliminaries on Semialgebraic Sets and CAD

A set \(S \subseteq \mathbb{R}^n\) is **semialgebraic** if it can be defined by a finite Boolean combination of polynomial inequalities and equations.

The **Cylindrical Algebraic Decomposition (CAD)** theorem asserts that:

- \(\mathbb{R}^n\) can be partitioned into finitely many **cells**.
- Each cell is **semialgebraic** and homeomorphic to an open cube \((0,1)^d\) for some \(d \leq n\).
- The cells are defined by conditions on the signs of a finite set of polynomials.

In particular, a cell is **open** in \(\mathbb{R}^n\) if and only if \(d = n\), and such a cell is described by a conjunction of **strict inequalities**:
\[
\{x \in \mathbb{R}^n : f_1(x) \sigma_1 0, \ldots, f_k(x) \sigma_k 0\}, \quad \sigma_i \in \{<, >\}.
\]

---

### Step 3. Proof of the Corollary for Open Semialgebraic Sets

Let \(S \subseteq \mathbb{R}^n\) be an **open** semialgebraic set.

1. By the CAD theorem, \(S\) is a finite union of cells:
   \[
   S = C_1 \cup C_2 \cup \cdots \cup C_m.
   \]

2. Since \(S\) is open, each cell \(C_j\) contained in \(S\) must be open in \(\mathbb{R}^n\). Hence, each \(C_j\) is of the form:
   \[
   C_j = \{x \in \mathbb{R}^n : f_1(x) \sigma_1 0, \ldots, f_k(x) \sigma_k 0\}, \quad \sigma_i \in \{<, >\}.
   \]

3. For each \(C_j\), we may replace any \(f_i\) with \(-f_i\) if necessary, so that all inequalities become strict **greater-than** inequalities. Thus, each \(C_j\) can be rewritten as:
   \[
   C_j = \{x \in \mathbb{R}^n : g_1(x) > 0, \ldots, g_k(x) > 0\}
   \]
   for some polynomials \(g_1, \ldots, g_k\).

4. Therefore, \(S\) is a finite union of sets of the desired form.

---

### Step 4. Conclusion

We have shown that every **open** semialgebraic subset of \(\mathbb{R}^n\) is a finite union of sets defined by strict polynomial inequalities of the form:
\[
\{x \in \mathbb{R}^n : f_1(x) > 0, \ldots, f_k(x) > 0\}.
\]

\[
\boxed{\text{The statement holds for open semialgebraic sets.}}
\]
\end{proof}
\end{document}

